Considera la funció $f(x)=x^2$.

\begin{apartats}

\apartat[1,25]{0,75}
Troba l'equació de la recta tangent a la gràfica de $f$ al punt d'abscissa $x=1$, i una primitiva de
$x^2-2x+1$.

\begin{solucio}
$f(1)=1$ i $f'(x)=2x$, $f'(1)=2$: la tangent és $y-1=2(x-1)$, és a dir, $y=2x-1$.\\
Com que $x^2-2x+1=(x-1)^2$, una primitiva és $\dfrac{(x-1)^3}3$.
\end{solucio}

\begin{tria}{area-tangent}
\itemtria{amb-eix-y}{1}{1,25}
Calcula l'àrea de la regió limitada per la gràfica de $f$, la recta tangent i l'eix $Y$.

\begin{solucio}
La paràbola és per sobre de la tangent, i totes dues es toquen a $x=1$. La regió va de $x=0$ a $x=1$:
\[
\int_0^1\bigl(x^2-(2x-1)\bigr)\,dx=\left[\frac{(x-1)^3}3\right]_0^1=0-\left(-\frac13\right)=\frac13
\]
unitats quadrades.
\end{solucio}

\itemtria{amb-eix-x}{1}{1,25}
Calcula l'àrea de la regió limitada per la gràfica de $f$, la recta tangent i l'eix $X$.

\begin{solucio}
La tangent talla l'eix $X$ a $x=\frac12$. La regió és la que hi ha sota la paràbola entre 0 i 1, menys
el triangle sota la tangent entre $\frac12$ i 1, de base $\frac12$ i altura 1:
\[
\int_0^1x^2\,dx-\frac12\cdot\frac12\cdot1=\frac13-\frac14=\frac1{12}
\]
unitats quadrades.
\end{solucio}
\end{tria}

\begin{nomesllarg}
\apartat{0,75}
Calcula l'àrea de la regió tancada per la paràbola $y=x^2$ i la recta $y=2x$.

\begin{solucio}
$x^2=2x$ dona $x=0$ i $x=2$, i entre aquests valors la recta és per sobre:
$\displaystyle\int_0^2(2x-x^2)\,dx=\left[x^2-\frac{x^3}3\right]_0^2=4-\frac83=\frac43$ unitats quadrades.
\end{solucio}
\end{nomesllarg}

\end{apartats}
